% ===========================================================================
\section{Implementation Details}
\label{app:details}
% ===========================================================================

\paragraph{Tokens.}
The vocabulary has 17 tokens: padding, start of sequence, \tok{Q}, \tok{=}, five symbols and eight function names.
The aligned layout uses four of the names (\tok{F1}, \tok{F2}, \tok{F3} for the tabulated functions and one for the induced function).
The material occupies 19 tokens after the start token: three rows of five values, then two key--value pairs for the induced function.
An exercise with a chain of length $d$ occupies $d+4$ tokens including its answer.

\paragraph{Model.}
A decoder-only transformer with pre-layer normalisation, width 64, four attention heads, an MLP expansion factor of four with GELU activations, learned absolute position embeddings added to the token embeddings, rotary position encoding inside the attention, and an output layer that is not tied to the input embedding.
A model described as ``$n$ layers $\times$ $r$ loops'' applies the same stack of $n$ layers $r$ times in sequence, with no input re-injection and no loop-index embedding, followed by one final layer normalisation.
The default model of experiments E1 and E2 is 2 layers $\times$ 2 loops (\NparamsMain{} parameters).

\paragraph{Training.}
AdamW with learning rate $2\times10^{-3}$, $\beta = (0.9, 0.98)$, weight decay $0.01$, linear warm-up over 200 steps followed by cosine decay to zero, gradient clipping at norm 1.0, and batches of 32 episodes with eight exercises each.
E1, E2 and E4 train for 5{,}000 steps (160{,}000 episodes); E3 trains for 8{,}000 steps (256{,}000 episodes).
The loss is the cross-entropy on answer tokens only.
In training, each of the three tabulated functions is an independent uniformly random map on the five symbols (except in one arm of E4, where it is a uniformly random permutation); the induced function is a cyclic shift by a uniformly random non-zero amount, shown through two examples with distinct random keys.

\paragraph{Evaluation.}
Each test consists of 4{,}000 episodes with a single exercise each, generated from a random stream that is independent of the training stream.
In tests on new textbooks the tabulated functions are uniformly random permutations.
In the G0 tests of E1 the textbooks are drawn from the finite training pool and are therefore arbitrary maps.
Curves during training use 500 episodes per test every 500 steps.

\paragraph{Guessing levels.}
On permutation textbooks the answer to a single look-up is uniform, so the guessing level of \skill{follow} is 20\%.
In a chain, the same function can occur more than once, and the answer then equals the input $x$ more often than $1/K$; for example $f(f(x)) = x$ whenever $x$ lies on a cycle of length one or two, which has probability $2/K$.
The best guesser that sees the exercise but not the tables answers $x$ for the chains where this is the likeliest value and anything else otherwise.
Enumerating all $120^3$ triples of permutations and all chains gives its accuracy exactly: \GuessTwo, \GuessThree, \GuessFour, \GuessFive{} and \GuessSix\% for chain lengths 2 to 6 (script \texttt{guess\_levels.py}).
For \skill{induce} the answer is never the input, so the guessing level is 25\%.
For \skill{mixed} at depth 2 the two functions are always different, and the level of 26.7\% derived in Section~\ref{sec:pilot} applies to a guesser that has induced $g$.

\paragraph{Compute.}
All \TotalRuns{} reported runs were executed on two virtual CPU cores without a GPU, using PyTorch 2.14 with one thread per run, and took \TotalCpuHours{} CPU-hours in total.

\paragraph{Release.}
The file \texttt{tinytextbook.py} contains the generator, the model, the training loop and the definitions of the four experiments; \texttt{make\_results.py} regenerates every table, figure and number in this paper from the JSON logs.
\texttt{guess\_levels.py} computes the guessing levels.
All three are available, together with the logs of every run, at \url{https://github.com/davitotty/Know-Less-Read-More}.

% ===========================================================================
\section{Per-Run Results}
\label{app:runs}
% ===========================================================================

Tables~\ref{tab:runs-e1}--\ref{tab:runs-e4} list the final accuracy (\%) of every run that enters the paper.
``Min.'' is the wall-clock training time in minutes.

\begin{table}[!htbp]
\centering\footnotesize
\begin{tabular}{@{}lrrrrrrrr@{}}
\toprule
Condition & Seed & Params & Steps & Min. & G0 fol. & G0 comp-2 & G1 fol. & G1 comp-2 \\
\midrule
M=1 & 0 & 113{,}297 & 5{,}000 & 6.9 & 100.0 & 100.0 & 20.7 & 21.3 \\
M=4 & 0 & 113{,}297 & 5{,}000 & 6.9 & 100.0 & 100.0 & 22.8 & 21.2 \\
M=16 & 0 & 113{,}297 & 5{,}000 & 7.0 & 100.0 & 100.0 & 31.4 & 21.1 \\
M=64 & 0 & 113{,}297 & 5{,}000 & 6.9 & 100.0 & 100.0 & 100.0 & 74.2 \\
M=64 & 1 & 113{,}297 & 5{,}000 & 6.8 & 100.0 & 99.6 & 81.0 & 28.4 \\
M=$\infty$ & 0 & 113{,}297 & 5{,}000 & 6.8 & 100.0 & 100.0 & 100.0 & 100.0 \\
M=$\infty$ & 1 & 113{,}297 & 5{,}000 & 6.6 & 100.0 & 100.0 & 100.0 & 100.0 \\
\bottomrule
\end{tabular}

\caption{E1, all runs.}
\label{tab:runs-e1}
\end{table}

\begin{table}[!htbp]
\centering\footnotesize
\begin{tabular}{@{}lrrrrrrrrr@{}}
\toprule
Condition & Seed & Params & Steps & Min. & follow & comp-2 & induce & mixed-2 & comp-3 \\
\midrule
full & 0 & 113{,}297 & 5{,}000 & 6.8 & 100.0 & 100.0 & 100.0 & 25.2 & 20.2 \\
full & 1 & 113{,}297 & 5{,}000 & 6.5 & 71.2 & 23.4 & 35.0 & 26.2 & 20.8 \\
full & 2 & 113{,}297 & 5{,}000 & 6.7 & 100.0 & 100.0 & 99.9 & 23.2 & 20.2 \\
full & 3 & 113{,}297 & 5{,}000 & 6.7 & 99.9 & 76.3 & 100.0 & 25.3 & 20.6 \\
no\_follow & 0 & 113{,}297 & 5{,}000 & 7.0 & 19.1 & 26.4 & 100.0 & 25.1 & 20.5 \\
no\_follow & 1 & 113{,}297 & 5{,}000 & 6.9 & 19.7 & 25.9 & 100.0 & 25.9 & 21.4 \\
no\_follow & 2 & 113{,}297 & 5{,}000 & 6.7 & 19.4 & 25.0 & 91.8 & 17.4 & 21.7 \\
no\_follow & 3 & 113{,}297 & 5{,}000 & 6.8 & 20.8 & 23.8 & 34.0 & 19.4 & 20.7 \\
no\_compose & 0 & 113{,}297 & 5{,}000 & 6.2 & 100.0 & 17.4 & 100.0 & 16.7 & 22.3 \\
no\_compose & 1 & 113{,}297 & 5{,}000 & 6.2 & 100.0 & 19.3 & 100.0 & 17.8 & 20.7 \\
no\_compose & 2 & 113{,}297 & 5{,}000 & 6.2 & 20.4 & 20.6 & 99.3 & 21.8 & 19.7 \\
no\_compose & 3 & 113{,}297 & 5{,}000 & 6.1 & 74.8 & 20.8 & 100.0 & 23.8 & 19.4 \\
no\_induce & 0 & 113{,}297 & 5{,}000 & 7.2 & 100.0 & 100.0 & 20.7 & 9.1 & 19.8 \\
no\_induce & 1 & 113{,}297 & 5{,}000 & 6.8 & 100.0 & 100.0 & 20.6 & 19.0 & 20.5 \\
no\_induce & 2 & 113{,}297 & 5{,}000 & 7.0 & 100.0 & 100.0 & 20.6 & 20.0 & 20.3 \\
no\_induce & 3 & 113{,}297 & 5{,}000 & 6.9 & 100.0 & 100.0 & 19.5 & 20.1 & 21.1 \\
\bottomrule
\end{tabular}

\caption{E2, all runs.}
\label{tab:runs-e2}
\end{table}

\begin{table}[!htbp]
\centering\footnotesize
\begin{tabular}{@{}lrrrrrrrrrr@{}}
\toprule
Condition & Seed & Params & Steps & Min. & len 1 & len 2 & len 3 & len 4 & len 5 & len 6 \\
\midrule
tied\_2x1 & 0 & 113{,}297 & 8{,}000 & 7.1 & 100.0 & 24.9 & 24.4 & 22.9 & 21.4 & 22.1 \\
tied\_2x2 & 0 & 113{,}297 & 8{,}000 & 13.6 & 100.0 & 100.0 & 85.6 & 22.0 & 21.5 & 22.6 \\
tied\_2x2 & 1 & 113{,}297 & 8{,}000 & 13.7 & 100.0 & 100.0 & 100.0 & 23.0 & 21.2 & 21.4 \\
tied\_2x2 & 2 & 113{,}297 & 8{,}000 & 13.2 & 100.0 & 100.0 & 30.0 & 25.1 & 23.8 & 21.6 \\
tied\_2x3 & 0 & 113{,}297 & 8{,}000 & 20.6 & 100.0 & 100.0 & 46.7 & 23.2 & 23.0 & 22.3 \\
untied\_4 & 0 & 213{,}265 & 8{,}000 & 13.7 & 100.0 & 76.4 & 30.1 & 22.7 & 20.6 & 21.2 \\
untied\_4 & 1 & 213{,}265 & 8{,}000 & 13.8 & 100.0 & 100.0 & 100.0 & 48.1 & 21.8 & 23.0 \\
untied\_4 & 2 & 213{,}265 & 8{,}000 & 13.4 & 100.0 & 91.5 & 34.9 & 23.1 & 22.4 & 21.2 \\
untied\_6 & 0 & 313{,}233 & 8{,}000 & 22.7 & 100.0 & 95.2 & 65.1 & 25.1 & 22.1 & 22.8 \\
\bottomrule
\end{tabular}

\caption{E3, all runs.}
\label{tab:runs-e3}
\end{table}

\begin{table}[!htbp]
\centering\footnotesize
\begin{tabular}{@{}lrrrrrr@{}}
\toprule
Condition & Seed & Params & Steps & Min. & perm. & maps \\
\midrule
train\_maps & 0 & 113{,}297 & 5{,}000 & 5.9 & 100.0 & 100.0 \\
train\_maps & 1 & 113{,}297 & 5{,}000 & 5.8 & 100.0 & 100.0 \\
train\_maps & 2 & 113{,}297 & 5{,}000 & 5.7 & 100.0 & 100.0 \\
train\_perms & 0 & 113{,}297 & 5{,}000 & 5.7 & 30.3 & 27.0 \\
train\_perms & 1 & 113{,}297 & 5{,}000 & 5.7 & 52.5 & 42.1 \\
train\_perms & 2 & 113{,}297 & 5{,}000 & 5.7 & 33.3 & 25.8 \\
\bottomrule
\end{tabular}

\caption{E4, all runs. ``perm.'' and ``maps'' are \skill{follow} accuracy on new textbooks that tabulate permutations and arbitrary maps.}
\label{tab:runs-e4}
\end{table}

\clearpage
% ===========================================================================
\section{Exploratory Runs}
\label{app:negative}
% ===========================================================================

The observations cited in Sections~\ref{sec:symmetry} and~\ref{sec:negative} come from runs made while the benchmark was being designed.
Their logs were not retained, with the exception of items~6 to~8; experiment E4 repeats the comparison of items~3 and~4 under controlled conditions.
They used a single seed each, were stopped early once the outcome seemed clear, and differ from one another in more than one factor, so they are reported as observations and not as controlled comparisons.
Unless noted, the models had four unshared layers of width 64 and all other settings were as in Appendix~\ref{app:details}.
Accuracies are on new permutation textbooks, where the guessing level of \skill{follow} is 20\%.

\begin{enumerate}
\item \textbf{Shuffled layout, permutations in training.}
With free-standing \tok{name key value} triples in random order, \skill{follow} stayed at chance for as long as we trained: 2{,}000 steps with exercises in postfix or prefix order and learned absolute positions; 1{,}500 steps with rotary encoding only; 3{,}000 steps when training on \skill{follow} alone with eight exercises per episode (one run at learning rate $10^{-3}$ with smaller textbooks mixed in, one at $3\times10^{-3}$ without).
\item \textbf{One tabulated function.}
Reducing the textbook to a single tabulated function plus the two examples of the induced one, with a two-layer model, left the training loss at $\ln 5$ for as long as we trained (3{,}000 to 3{,}800 steps), with and without an added next-token loss on all positions.
A model of the same size fitted one fixed batch of 32 full-size episodes to a loss below $0.01$ in 200 steps, so the failure is one of finding the general solution and not of capacity or of a faulty pipeline.
\item \textbf{Aligned layout, permutations in training.}
\skill{follow} was at chance at every checkpoint up to step 3{,}500 of a 4{,}000-step schedule at learning rate $10^{-3}$, while \skill{induce}, whose two examples sit at fixed positions, reached 100\% by step 1{,}500.
\item \textbf{Aligned layout, arbitrary maps in training.}
With no other change, \skill{follow} went from 32\% at step 2{,}000 to 99.8\% at step 2{,}500; at learning rate $2\times10^{-3}$ it reached 98\% by step 1{,}500.
\item \textbf{Shuffled layout, arbitrary maps in training.}
\skill{follow} was still at chance when the run was stopped at step 2{,}500.
\item \textbf{Four unshared layers on the E2 training mix.}
With \skill{follow}, \skill{induce} and \skill{compose} at depths 2 and 3 in the training mix and learning rate $2\times10^{-3}$, \skill{follow} and \skill{induce} reached 100\%, but \skill{compose} at depth 2 had only reached 29\%, barely above its guessing level, when the run was stopped at step 7{,}000 of 10{,}000.
This observation motivated both the looped default model and experiment E3; E3 then showed that the difference between looped and unshared models is not consistent across seeds, so the observation should not be read as evidence about architecture.
\item \textbf{Full curriculum with a 6{,}000-step schedule.}
Before the budget was fixed, the full curriculum was run once with the final default model and a schedule of 6{,}000 steps (seed 0).
It ended with \skill{follow} and \skill{induce} at 100\%, \skill{compose} at depth 2 at 82.7\% and \skill{mixed} at 21.7\%.
The run is not part of Table~\ref{tab:e2} because its budget differs from that of the reported runs; it is one more instance of a full-curriculum run that did not solve \skill{compose}.
\item \textbf{Three loops with a 10{,}000-step schedule.}
The three-loop model of E3, trained with a 10{,}000-step schedule (seed 0), was at 95.2\% at chain length 3 by step 6{,}500 and still at the guessing level at length 4 when it was stopped at step 7{,}500.
\end{enumerate}

